\section{Related work}
\label{sec:related}

\paragraph{Contrastive learning with and without labels}
\simclr \citep{chen2020simclr} contrasts two views of an image against the other images of a large batch;
MoCo replaces the large batch by a queue of past embeddings \citep{he2020moco,chen2020mocov2}; BYOL and
SimSiam remove the negatives \citep{grill2020byol,chen2021simsiam}. \supcon \citep{khosla2020supcon} treats all
images of a class as positives and, with all labels and a ResNet-50, reports 96.0\% on CIFAR-10 against
93.6\% for \simclr and 95.0\% for cross-entropy. That comparison uses every label; this paper asks what
happens when only a small fraction is available.

\paragraph{Labels inside the loss}
SuNCEt \citep{assran2020suncet} adds a supervised contrastive term on the labeled images to \simclr and
shows that, with 10\% of the ImageNet labels, it matches the semi-supervised accuracy of \simclr with less
than half of the pretraining compute. \citet{inoue2020semi} write supervised and unsupervised contrastive
learning as one generalised loss, and \citet{islam2021broad} study joint \simclr and \supcon objectives with
all labels. PAWS \citep{assran2021paws} and S$^4$L \citep{zhai2019s4l} combine self-supervised and supervised
signals in other ways. A second family starts from pseudo-labeling and adds a contrastive regulariser:
CoMatch \citep{li2021comatch}, SimMatch \citep{zheng2022simmatch}, class-aware contrastive semi-supervised
learning \citep{yang2022classaware} and the clustering-based method of \citet{fini2023semi} all combine labels
and instance-level learning successfully, and USB \citep{wang2022usb} standardises their evaluation. Our
\union arm is the simplest member of the first family; we use it as a controlled probe, not as a
competitor to these methods.

\paragraph{Labels after pretraining}
SimCLR~v2 \citep{chen2020simclrv2} pretrains without labels, fine-tunes on the few labels and then distils
the fine-tuned network on unlabeled data; CPC~v2 \citep{henaff2020data} showed earlier that self-supervised
features allow training with far fewer labels. \citet{zoph2020rethinking} compare pretraining with
self-training across label fractions, and \citet{newell2020useful} measure how the benefit of
self-supervised pretraining shrinks as labels grow. These works establish that the pretrain-then-use-labels
recipe is strong. They do not compare it, at matched compute, with putting the same labels inside the
contrastive loss.

\paragraph{Class collapse and transfer}
At the end of supervised training the features of a class collapse to their mean
\citep{papyan2020neural}; \citet{galanti2022role} measure this with the class-distance normalised variance
(CDNV), and \citet{hui2022limitations} show that collapse on training data need not carry over to test data.
For \supcon, \citet{graf2021dissecting} analyse the collapsed optimum, \citet{chen2022perfectly} show that
class collapse hurts transfer and robustness and counter it with a class-conditional instance term, and
\citet{xue2023which} explain collapse through the simplicity bias of gradient descent.
\citet{kornblith2021better} find that losses with tighter class separation give features that transfer worse,
and \citet{zhao2021what} that the within-class invariance of supervised pretraining discards information
that transfer needs. Our diagnostic applies these ideas to the label-scarce case, where the question is
whether tightness on the \emph{labeled} images reaches the \emph{unlabeled} ones.

\paragraph{Coverage and selection of small sets}
Core-set selection picks points that cover the data \citep{sener2018coreset}. For small budgets, typical
examples beat atypical ones \citep{hacohen2022active}, coverage of the embedding space is the right
selection objective \citep{yehuda2022active}, and data pruning should keep easy examples when data are
scarce \citep{sorscher2022beyond,zheng2023coverage}. We do not propose a selection method; we use pools of
controlled coverage to test whether these findings also explain how a labeled pool limits \supcon.

\paragraph{Budgets and evaluation}
\citet{li2020budgeted} argue for comparing methods at a fixed iteration budget, and
\citet{dehghani2022efficiency} show that a single cost indicator can mislead.
\citet{oliver2018realistic} criticise semi-supervised evaluations that tune on validation sets larger than
the labeled set, and \citet{newell2020useful} note that linear evaluation does not predict fine-tuning
performance. \citet{kumar2022lpft} show that fine-tuning can distort pretrained features and propose to
start from a linear probe. For reference, the solo-learn library reports CIFAR-10 linear accuracies between
88.7\% and 92.9\% for eleven self-supervised methods trained much longer than our reference budget
\citep[Table~1]{dacosta2022sololearn}.

\paragraph{What is new here}
\Cref{tab:novelty} places each of our findings next to the closest earlier evidence. The individual
phenomena have precedents. What was missing is (i) a comparison in which labels and images are varied
separately at equal compute, (ii) a demonstration that the choice between the all-label probe and a
label-budget protocol changes the answer, and (iii) a measurement of tightness on labeled, unlabeled and
unseen images for label-scarce pretraining.

\begin{table}[t]
\centering
\small
\caption{Our findings next to the closest prior evidence.}
\label{tab:novelty}
\begin{adjustbox}{max width=\linewidth}
\begin{tabular}{p{0.24\linewidth} p{0.34\linewidth} p{0.38\linewidth}}
\toprule
finding & closest prior evidence & what this paper adds \\
\midrule
Labels in the contrastive loss speed up pretraining &
SuNCEt reaches the accuracy of \simclr with less than half the compute at 10\% ImageNet labels
\citep{assran2020suncet} &
The speed-up is isolated from the choice of images in a $2\times2$ design at equal steps, and followed to
long budgets, where it turns into a deficit \\
\addlinespace
Labels are best used after label-free pretraining &
SimCLR~v2 \citep{chen2020simclrv2}; pretraining versus self-training \citep{zoph2020rethinking} &
A direct, compute-matched comparison with the same labels inside a \supcon loss, including two-stage and
semi-supervised alternatives \\
\addlinespace
The evaluation protocol can change conclusions &
Unrealistic validation sets in semi-supervised learning \citep{oliver2018realistic}; linear evaluation does
not predict fine-tuning \citep{newell2020useful} &
How much the protocol and the split of the labels change the comparison of the same encoders: default, selectable and oracle fine-tuning recipes, and labels held out for selection at two budgets \\
\addlinespace
Tight classes on training data need not generalise &
Collapse on train but not on test \citep{hui2022limitations}; class collapse in \supcon hurts transfer
\citep{chen2022perfectly,kornblith2021better,xue2023which} &
A three-way split (labeled, unlabeled, test) for label-scarce pretraining, with a cross-entropy control \\
\addlinespace
Coverage matters for small pools &
Low-budget active learning and data pruning
\citep{hacohen2022active,yehuda2022active,sorscher2022beyond} &
An intervention on the labeled pool of \supcon at fixed size; coverage explains part of the effect \\
\bottomrule
\end{tabular}
\end{adjustbox}
\end{table}
